\begin{proposition}[5.13.7]
\label{IV.2.s5.13.7-fr}
Supongamos que los $A_\alpha$ son regulares y que, para
$\alpha\leq\beta$, $\varphi_{\beta\alpha}$ hace de $A_\beta$ un
$A_\alpha$-módulo plano; supongamos además que
$A=\varinjlim A_\alpha$ es noetheriano. Entonces $A$ es regular.
\end{proposition}

Por definición, basta probar que, para todo ideal primo $\mathfrak p$ de $A$,
el anillo local noetheriano $A_{\mathfrak p}$ es regular; ahora bien, si
$\mathfrak p_\alpha=\varphi_\alpha^{-1}(\mathfrak p)$, se tiene
$\mathfrak p=\varinjlim\mathfrak p_\alpha$ y
$A_{\mathfrak p}=\varinjlim((A_\alpha)_{\mathfrak p_\alpha})$ por (5.13.3);
como, por hipótesis, los $(A_\alpha)_{\mathfrak p_\alpha}$ son regulares y
$(A_\beta)_{\mathfrak p_\beta}$ es un
$(A_\alpha)_{\mathfrak p_\alpha}$-módulo plano para
$\alpha\leq\beta$ $(\mathbf{0}_{\mathrm I}, 6.3.2)$, se ve que la cuestión se reduce al
caso en que los $A_\alpha$ y $A$ son además anillos \emph{locales}. Para probar que
$A$ es regular, basta entonces mostrar que, si $M$, $N$ son dos $A$-módulos de
tipo finito, existe $i_0$ tal que, para $i\geq i_0$,

\oldpage[IV-2]{134}

se tiene $\operatorname{Tor}^A_i(M,N)=0$ (\textbf 0, 17.3.1 y 17.2.6). Pero se tiene
el lema siguiente (en el que los $A_\alpha$ son de nuevo cualesquiera):

\begin{lemma}[5.13.7.1]
\label{IV.2.s5.13.7.1-fr}
Para todo $A$-módulo $M$ de presentación finita, existe un índice $\alpha$ y
un $A_\alpha$-módulo de presentación finita $M_\alpha$ tal que $M$ sea
isomorfo a $M_\alpha\otimes_{A_\alpha}A$.
\end{lemma}

En efecto, se puede suponer que $M=A^r/P$, donde $P$ es un sub-$A$-módulo de
tipo finito de $A^r$. Como $A^r=\varinjlim(A_\alpha^r)$ y $P$ es de tipo
finito, existe un $\alpha$ y un sub-$A_\alpha$-módulo de tipo finito
$P_\alpha$ de $A_\alpha^r$ tal que $P$ sea la imagen canónica de
$P_\alpha\otimes_{A_\alpha}A$; la exactitud por la derecha del producto
tensorial muestra, pues, que $M$ es isomorfo a
$(A_\alpha^r/P_\alpha)\otimes_{A_\alpha}A$.

Establecido este lema, $M$ y $N$ son de presentación finita, puesto que $A$ es
noetheriano; escribamos, por tanto,
$M=M_\alpha\otimes_{A_\alpha}A$,
$N=N_\alpha\otimes_{A_\alpha}A$; en virtud de la hipótesis, $A$ es un
$A_\alpha$-módulo plano $(\mathbf{0}_{\mathrm I}, 6.1.2)$, por lo que se tiene
$\operatorname{Tor}^A_i(M,N)
=\operatorname{Tor}^{A_\alpha}_i(M_\alpha,N_\alpha)
 \otimes_{A_\alpha}A$
salvo isomorfismo (\textbf{III}, 6.3.9); como, por hipótesis, $A_\alpha$ es regular,
existe $i_0$ tal que
$\operatorname{Tor}^{A_\alpha}_i(M_\alpha,N_\alpha)=0$ para $i\geq i_0$.
C.Q.F.D.

\begin{corollary}[5.13.8]
\label{IV.2.s5.13.8-fr}
Supongamos que los $A_\alpha$ son noetherianos y verifican la propiedad
$(R_k)$ para un entero $k\geq0$. Supongamos además que, para
$\alpha\leq\beta$, $\varphi_{\beta\alpha}$ hace de $A_\beta$ un
$A_\alpha$-módulo plano, que $A=\varinjlim A_\alpha$ es noetheriano y que,
para todo $\alpha$, ${}^a\varphi_\alpha$ es un homeomorfismo de
$\operatorname{Spec}(A)$ sobre $\operatorname{Spec}(A_\alpha)$. Entonces $A$
verifica $(R_k)$.
\end{corollary}

Sea, en efecto, $\mathfrak p$ un ideal primo de $A$ y, para todo $\alpha$,
pongamos $\mathfrak p_\alpha=\varphi_\alpha^{-1}(\mathfrak p)$, de modo que
$\mathfrak p=\varinjlim\mathfrak p_\alpha$ y
$A_{\mathfrak p}=\varinjlim(A_\alpha)_{\mathfrak p_\alpha}$ (5.13.3). De
(\textbf I, 2.4.2) y de la hipótesis resulta que el morfismo
$\operatorname{Spec}(A_{\mathfrak p})
\longrightarrow\operatorname{Spec}((A_\alpha)_{\mathfrak p_\alpha})$
es un homeomorfismo sobreyectivo para todo $\alpha$; por tanto, se tiene
$\dim(A_{\mathfrak p})=\dim((A_\alpha)_{\mathfrak p_\alpha})$ para todo
$\alpha$. Supongamos entonces que $\dim(A_{\mathfrak p})\leq k$; se tiene,
por tanto, $\dim((A_\alpha)_{\mathfrak p_\alpha})\leq k$, y la hipótesis
implica que $(A_\alpha)_{\mathfrak p_\alpha}$ es un anillo regular; como,
además, para $\alpha\leq\beta$,
$(A_\beta)_{\mathfrak p_\beta}$ es un
$(A_\alpha)_{\mathfrak p_\alpha}$-módulo plano
$(\mathbf{0}_{\mathrm I}, 6.3.2)$, se concluye de (5.13.7) que
$A_{\mathfrak p}$ es regular, lo que demuestra el corolario.
